\section*{Lectures on Lie Groups and Lie Algebras – 8,9,10}

\section*{Lie Groups}

\textbf{Definition 1.}
A \emph{Lie group} is a group $G$ equipped with a smooth manifold structure
compatible with the group operations. This means:

\begin{enumerate}
\item The multiplication map
\[
m : G \times G \to G, \quad m(g,h)=gh
\]
is smooth.

\item The inversion map
\[
s : G \to G, \quad s(g)=g^{-1}
\]
is smooth.
\end{enumerate}

Each element $g\in G$ defines diffeomorphisms
\[
L_g(h)=gh, \qquad R_g(h)=hg.
\]

\textbf{Examples.}

\begin{enumerate}
\item Any discrete (in particular finite) group;
\item $\mathbb{R}$ under addition;
\item $\mathbb{R}_+$ under multiplication;
\item $S^1 = \{z\in\mathbb{C} : |z|=1\}$;
\item $GL_n(\mathbb{R})$, $GL_n(\mathbb{C})$;
\item $SL_n(\mathbb{R})$, $SL_n(\mathbb{C})$.
\end{enumerate}

\textbf{Definition 2.}
A Lie group homomorphism is a smooth group homomorphism.
A bijective one is an isomorphism.

Examples:
\[
t \mapsto e^t : \mathbb{R}\to\mathbb{R}_+
\]
\[
t \mapsto e^{2\pi i t} : \mathbb{R}\to S^1
\]
\[
\det : GL_n(\mathbb{R}) \to \mathbb{R}^\times.
\]

\textbf{Definition 3.}
A Lie subgroup is a closed smooth submanifold closed under multiplication.

Examples:
\[
SL_n \subset GL_n, \quad N_+(n), \quad B_+(n), \quad O_n(\mathbb{R}), \quad SO_n(\mathbb{R}).
\]

\textbf{Proposition 1.}
If $H\subset G$ is a Lie subgroup, then $G/H$ carries a natural manifold structure.
If $H$ is normal, then $G/H$ is a Lie group.

\textbf{Proposition 2. (Homomorphism Theorem)}
If $\varphi:G_1\to G_2$ is surjective, then
\[
G_2 \cong G_1/\ker\varphi.
\]

\section*{Classification of 1-Dimensional Lie Groups}

\textbf{Proposition 3.}
Every connected 1-dimensional Lie group is isomorphic to
\[
\mathbb{R} \quad \text{or} \quad S^1.
\]

\section*{The Lie Functor}

\textbf{Definition 5.}
A vector field $\xi$ is left-invariant if
\[
L_{g*}\xi=\xi.
\]

\textbf{Proposition 4.}
A left-invariant vector field is determined by its value at the identity.

The Lie algebra of $G$ is
\[
\mathfrak{g} = \mathrm{Lie}\,G = T_eG.
\]

\textbf{Definition 6.}
If
\[
m(x,y)=x+y+B(x,y)+o(x,y),
\]
then
\[
[x,y]=B(x,y)-B(y,x).
\]

Equivalent definitions:
\[
[x,y]=\frac{d}{dt}\frac{d}{ds}
g_x(t)g_y(s)g_x(t)^{-1}g_y(s)^{-1}
\]

or

\[
[x,y]=[\xi_x,\xi_y](e).
\]

\textbf{Adjoint Representation}

\[
\mathrm{Ad}(g)=d_e(L_gR_{g^{-1}}).
\]

Examples:
\[
\mathrm{Lie}\,\mathbb{R}=\mathbb{R},
\quad
\mathrm{Lie}\,GL_n=\mathfrak{gl}_n,
\quad
\mathrm{Lie}\,SO_n=\mathfrak{so}_n.
\]

\textbf{Proposition 5.}
If $\varphi:G_1\to G_2$, then
\[
d_e\varphi:\mathrm{Lie}\,G_1\to \mathrm{Lie}\,G_2
\]
is a Lie algebra homomorphism.

\section*{Representations of 1-Dimensional Lie Groups}

\textbf{Proposition 6.}
Every finite-dimensional representation of $\mathbb{R}$ has the form
\[
\rho(t)=\exp(tA).
\]

\textbf{Proposition 7.}
Every finite-dimensional representation of $S^1$
is a direct sum of 1-dimensional representations
\[
\rho(z)=z^k, \quad k\in\mathbb{Z}.
\]

\section*{$SU(2)$ and $SO(3)$}

\textbf{Definition 10.}
\[
SU(2)=\{A\in GL_2(\mathbb{C}) : A^\dagger=A^{-1}, \det A=1\}.
\]

\textbf{Proposition 8.}
As a manifold,
\[
SU(2)\cong S^3.
\]

\textbf{Proposition 9.}
\[
\mathrm{Ad}:SU(2)\to SO(3)
\]
is a double covering.

Kernel:
\[
\{\pm I\}.
\]

\section*{Connectedness and Universal Covering}

\textbf{Proposition 10.}
The identity component $G^0$ is a Lie subgroup.

\textbf{Theorem 1.}
Every connected Lie group has a unique simply connected universal cover.

Examples:
\[
\mathbb{R}\to S^1,
\quad
SU(2)\to SO(3).
\]

\section*{Fundamental Groups}

\textbf{Proposition 11.}
\[
\pi_1(SO_n(\mathbb{R}))=\mathbb{Z}/2\mathbb{Z}, \quad n\ge3.
\]

\textbf{Proposition 12.}
\[
\pi_1(SL_2(\mathbb{R}))=\mathbb{Z}.
\]

\section*{Haar Measure and Complete Reducibility}

\textbf{Proposition 13.}
Every Lie group admits a left-invariant measure (Haar measure),
unique up to scalar.

Examples:
\[
dt \text{ on } \mathbb{R}, \quad d\varphi \text{ on } S^1.
\]

\textbf{Theorem 2.}
Every finite-dimensional representation of a compact Lie group
is completely reducible.

Proof: Average a projector:
\[
P=\int_G \rho(g)p\rho(g^{-1})\,dg.
\]

Then $P$ is an invariant projection and
\[
V=W\oplus\ker P.
\]